\documentclass[11pt, a4paper]{article}
\usepackage[top=2.5cm, bottom=2.5cm, left=2.8cm, right=2.8cm]{geometry}
\usepackage{palatino}
\usepackage{setspace}
\singlespacing
\usepackage{graphicx}
\usepackage{caption}
\usepackage{float}
\usepackage{needspace}
\usepackage{amsmath}
\usepackage{booktabs}
\usepackage{microtype}
\usepackage{fancyhdr}
\usepackage{titlesec}
\usepackage[numbers]{natbib}
\usepackage[hidelinks]{hyperref}
\usepackage{gensymb}

\captionsetup{
  font=small,
  labelfont=bf,
  format=plain,
  justification=justified,
  skip=4pt
}

\setlength{\parskip}{6pt}
\setlength{\parindent}{0pt}
\setlength{\intextsep}{10pt}
\setlength{\textfloatsep}{10pt}
\setlength{\floatsep}{8pt}

\titleformat{\section}{\large\bfseries}{}{0em}{}
\titlespacing{\section}{0pt}{14pt}{4pt}

\pagestyle{fancy}
\fancyhf{}
\lhead{\small\textit{Ravulapalli --- Landmark Symmetry and Cognitive Map Formation}}
\rhead{\small\thepage}
\renewcommand{\headrulewidth}{0.4pt}

\fancypagestyle{plain}{
  \fancyhf{}
  \rhead{\small\thepage}
  \renewcommand{\headrulewidth}{0pt}
}

\graphicspath{{images/}{figures/}}
\usepackage{amssymb}

\begin{document}

\begin{center}
  {\LARGE\bfseries Landmark Symmetry and Cognitive Map Formation\\[6pt]
  in Predictive Recurrent Neural Networks}\\[14pt]
  {\large Manas Venkata Sai Ravulapalli}\\[4pt]
  {\normalsize Independent Study Module --- Computational Neuroscience}
\end{center}


\thispagestyle{plain}

\begin{abstract}
Predictive learning has been proposed as a general mechanism for forming
cognitive maps in the hippocampus: a recurrent neural network trained to
predict future observations from the statistics of an animal's movement
will develop place-cell-like spatial representations~\citep{levenstein2024}.
A natural question is how this framework behaves when the environment
provides ambiguous sensory information. We extend the Levenstein et al.
pRNN framework to ask: does rotational landmark symmetry cause map
formation to fail?

We train predictive RNNs in arenas with C1, C2, and C4 rotational
landmark symmetry and evaluate the resulting representations using four
novel metrics designed specifically to detect the forms of failure that
standard spatial RSA cannot resolve: Permutation-Aware Alignment (PAA),
Rotational Autocorrelation (RA), C2 Contrast, and a Symmetry Collapse
Index (SCI). We find that map formation does not fail under C4 symmetry.
Instead, a stable spatial map forms, but is geometrically distorted: units
develop rate maps that are similar at rotationally equivalent positions,
and decoding precision degrades monotonically with symmetry order
(RA: S1\,=\,0.097, S2\,=\,0.114, S4\,=\,0.217). Global orientational degeneracy is prevented by the heading direction
(HD) input: PAA gain remains near zero for all conditions
(S4\,$<$\,0.05). Removing HD input causes map formation to collapse
(sRSA: 0.615 $\to$ 0.260), confirming
that HD is necessary for global stability. Together, these results establish
that predictive learning is robust to landmark symmetry in terms of whether
a map forms, but sensitive to it in terms of the geometric quality of the
formed representation.
\end{abstract}



\section*{Introduction}

Levenstein et al.~(2024) showed that a recurrent network trained on sequential predictive loss spontaneously develops spatially structured representations---but validated this exclusively in an L-shaped arena chosen to ensure every position has a unique visual fingerprint. Most rodent electrophysiology experiments use square arenas. Square arenas are rotationally symmetric. The condition Levenstein's validation depends on---global observational discriminability---does not hold there. We ask what the Levenstein framework actually does under these conditions. Not what the hippocampus does. What the model does. The answer has direct consequences for how the framework should be used, what its representations are grounded in, and what any extension to more naturalistic environments should expect to find.

The head-direction (HD) system tracks global orientation through path integration while landmark anchoring corrects drift via visual cues~\cite{taube1990, knierim1995}, but symmetric landmarks force these two mechanisms into direct conflict: the same landmark view appears at multiple rotations, so the HD state and the landmark signal become discordant. The double-rotation literature documents that place cells increase their number of firing fields in symmetric arenas~\cite{muller1987, knierim1995}, with multi-field fraction scaling with symmetry order, and these multi-field patterns emerge spontaneously in the hippocampus without explicit instruction~\cite{leutgeb2005}.

This paper asks whether landmark symmetry causes global map degeneracy (where different seeds learn incompatible global orientations), local representational folding (where symmetric positions collapse within a single run), or both. A standard metric (seed-to-seed RSA) cannot answer this question because row permutations preserve pairwise distances, making it blind to both global rotations and local folding. We introduce three metrics that separate identifiability from entanglement: PAA and RA for global vs local, C2 Contrast for group structure, and use them to test predictive networks trained on arenas with graded rotational symmetry. We show that head-direction input prevents global degeneracy while local precision degrades monotonically with symmetry order, and that this partial-folding pattern explains multi-field place cell emergence, predicts specific lesion deficits, and uncovers why sRSA fails under observation aliasing.

The theoretical link between predictive learning and spatial representation
comes from the successor representation (SR)~\citep{stachenfeld2017}. The
SR encodes each state as the discounted expected future occupancy of all
other states: $M(s,s') = \mathbb{E}\left[\sum_{t=0}^{\infty} \gamma^t
\mathbf{1}(s_t = s') \mid s_0 = s\right]$. Under a random policy in an
open environment, the SR matrix has a structure similar to a Gaussian
centered at each state, producing place-cell-like activity
profiles~\citep{stachenfeld2017}. Critically, the SR geometry is shaped
by the transition structure of the environment, not just its physical
geometry. This means that environments with rotational symmetry will have
SR matrices in which multiple distinct spatial positions are predicted to
share the same future state distribution --- the SR entries for these
positions will be identical or nearly so. Whether a predictive RNN can
resolve this ambiguity, and if not, in what way its representation fails,
is the central question of this work.

Here we extend the pRNN framework of~\citet{levenstein2024} to
systematically ask: how does rotational landmark symmetry affect the
quality of the learned spatial representation? We address
this with two interleaved experiments. First, a symmetry sweep across
three arena conditions (C1, C2, C4) with matched object density, analyzed
with four metrics designed to detect forms of representational failure that
spatial RSA cannot resolve. Second, an HD ablation experiment that removes
or degrades the heading direction input to test whether HD is necessary for
global representational stability. Across both experiments, we find that the pRNN framework is more robust to environmental ambiguity --- but this robustness has clear geometric limits that depend on both the structure of the environment and the quality of the agent's sensory signal.



\section*{Methods}

The arena is an $18 \times 18$ MiniGrid with 324 positions. All conditions are matched on observational discriminability:
\[
\mathrm{ODI}=1-\frac{H(X\mid O)}{H(X)}=0.336.
\]
S4 has C4-equivalent observations under $90^{\circ}$ rotation ($H_2=4.0$), S2 has C2-equivalent observations under $180^{\circ}$ rotation ($H_2=2.0$), and S1 has no designed equivalence ($H_2=1.0$). The number of seeds is S4 $n=7$, S2 $n=5$, S1 $n=5$.

The model is a NormReLU pRNN with $N=500$ units. At time $t$, the hidden state updates as:
\[
h_t = \mathrm{NormReLU}(W_h h_{t-1} + W_o o_t + W_a a_t + W_d d_t),
\]
where $o_t$ is the egocentric observation, $a_t$ is the action, and $d_t$ is a 5-dimensional SpeedHD vector (speed + 4-way one-hot heading direction). Training minimises multi-step prediction loss for $k=5$ steps:
\[
\mathcal{L} = \sum_{j=1}^{k}\|\hat{o}_{t+j} - o_{t+j}\|_2^2,
\]
over 80,000 training steps. A total of 17 seeds were run across all conditions (S4 $n=7$, S2 $n=5$, S1 $n=5$).

\begin{figure}[tbp]
\centering
\begin{minipage}{0.48\textwidth}
  \centering
  \includegraphics[width=\textwidth]{images/fig_arenas_overview.png}
  \caption{Arena layouts by symmetry. S4 has $C_4$ symmetry; S2 has $C_2$; S1 is asymmetric.}
  \label{fig:arena}
\end{minipage}
\hfill
\begin{minipage}{0.48\textwidth}
  \centering
  \includegraphics[width=\textwidth]{images/fig_srsa_invariance.png}
  \caption{Why sRSA is invariant to permutations. Both maps yield identical pairwise distances.}
  \label{fig:srsa}
\end{minipage}
\end{figure}

\noindent Three purpose-built metrics address this blindness. PAA detects when rotated alignment between seeds beats identity alignment, revealing orientational divergence. RA measures whether units respond identically at rotationally equivalent positions, capturing unit-level aliasing. C2 Contrast tests whether the population code encodes group structure. PAA, RA, and C2 Contrast separate identifiability, entanglement, and population structure---what sRSA cannot do.

\section*{sRSA cannot detect the failure modes this paper measures}

Seed-to-seed RSA correlates pairwise neural distances with pairwise spatial distances. However, this metric has a critical vulnerability: row permutations of the neural population vector preserve all pairwise distances, so rotated maps yield identical sRSA values even though they represent opposite global orientations.

Figure \ref{fig:srsa} demonstrates this blindness: two fundamentally different failure modes---global rotation versus local folding---yield identical sRSA despite producing qualitatively different neural codes. This blindness matters empirically: sRSA is near-equal for S2 and S4 (~0.72 vs ~0.67), yet RA differs by nearly 2× (0.114 vs 0.217). A researcher relying on sRSA alone would conclude that C2 and C4 symmetry impose equivalent representational costs. They do not.

\section*{Three possible outcomes frame every result that follows}

The space of possible outcomes is structured by two independent questions: Does global orientation remain stable across seeds? Does local precision stay high? This produces three mechanistically distinct failure modes: (1) global degeneracy and local collapse, (2) global stability but local folding (the case we observe), or (3) full robustness.

\begin{figure}[tbp]
\centering
\begin{minipage}{0.48\textwidth}
  \centering
  \includegraphics[width=\textwidth]{images/fig_failure_modes.png}
  \caption{Mechanistically distinct outcomes. Left: collapse. Center: partial folding (observed). Right: null case.}
  \label{fig:outcomes}
\end{minipage}
\hfill
\begin{minipage}{0.48\textwidth}
  \centering
  \includegraphics[width=\textwidth]{figures/fig_metric_summary.pdf}
  \caption{Final metric comparison. PAA (global), RA (aliasing), SCI (compression), and C2 Contrast (grouping).}
  \label{fig:dash}
\end{minipage}
\end{figure}

\noindent These four metrics reveal the pattern: global structure is stable (PAA $< 0.05$), but local precision degrades with symmetry (RA increases with $|C_n|$), and the population code encodes the observation symmetry group itself (C2 Contrast negative in S2, positive in S4).

\section*{HD input prevents global orientational degeneracy}

Global map degeneracy did not occur. Cross-seed RSA alignment is S4 $0.977 \pm 0.002$, S2 $0.975 \pm 0.003$, S1 $0.974 \pm 0.001$. PAA gain is S4 $0.0015 \pm 0.0002$, S2 $0.0021 \pm 0.0004$, S1 $0.0000 \pm 0.0001$ with $p = 0.786$ and $r = -0.200$---indistinguishable across conditions and all below 0.05.

\begin{figure}[tbp]
\centering
\begin{minipage}{0.48\textwidth}
  \centering
  \includegraphics[width=\textwidth]{figures/fig_rdm_comparison.pdf}
  \caption{Cross-seed RDM alignment. S4, S2, and S1 all show strong RSA alignment across seeds.}
  \label{fig:rdm}
\end{minipage}
\hfill
\begin{minipage}{0.48\textwidth}
  \centering
  \includegraphics[width=\textwidth]{figures/fig_metric_trajectories.pdf}
  \caption{Distance-Topology Gap ($\Delta_{\text{TG}}$). Gap widens in S4, indicating path-dependence.}
  \label{fig:dtg}
\end{minipage}
\end{figure}

\noindent Even with $C_4$-symmetric observations, the input stream itself is not symmetric: moving north from position $(r,c)$ and from its $90^{\circ}$ equivalent $(c, 17-r)$ yields different head-direction-tagged transitions at different absolute orientations. These transitions remain globally consistent across any arena trajectory, so the network must track global heading precisely to minimize prediction loss. Global orientation therefore locks in deterministically at training initialization. This closes the first failure mode: global degeneracy is ruled out.

\section*{Unit-level aliasing scales monotonically with symmetry order}

Head-direction input cannot prevent local aliasing costs. Unit-level rotational autocorrelation (RA) was S4 0.217, S2 0.114, S1 0.097 --- scaling directly with symmetry group order (4-fold, 2-fold, 1-fold respectively). S4 versus S1 yields $p = 0.036$ and $r = -1.0$. Spatially tuned units (EVS $> 0.1$) comprise S4 1.8\%, S2 3.7\%, S1 4.7\% of the population.

\begin{figure}[tbp]
\centering
\includegraphics[width=0.85\textwidth]{figures/fig_rate_maps.pdf}
\caption{Example place cell rate maps for the six highest-EVS units per condition (seed\_00). \textbf{S1:} compact, idiosyncratic single fields typical of healthy place cells. \textbf{S2:} bilateral or 4-fold symmetric fields --- the same firing pattern duplicated at rotationally equivalent positions. \textbf{S4:} cross-shaped and stripe-like patterns emerge; some units are nearly silent; fields are substantially broader than S1. These patterns are the unit-level signature of partial folding: the network encodes within-quadrant position but cannot disambiguate which quadrant.}
\label{fig:tune}
\end{figure}

\begin{figure}[tbp]
\centering
\includegraphics[width=0.65\textwidth]{figures/fig_field_size.pdf}
\caption{Field width distributions by condition. Mean field width increases monotonically from S1 to S4, with a subpopulation of S4 units developing fields spanning the full arena width.}
\label{fig:fieldsize}
\end{figure}

\noindent High-RA units in S4 produce four distinct peaks across quadrants---they encode within-quadrant position but cannot disambiguate which quadrant. Symmetric observations remove the prediction gradient that would force quadrant differentiation, so aliasing accumulates with symmetry order. The field width distributions (Figure~\ref{fig:fieldsize}) confirm this: mean field width increases monotonically from S1 to S4, with a subpopulation of S4 units developing fields spanning the full arena width.

\section*{Decoding degrades while population collapse stays partial}

Population compression stays graded despite unit-level aliasing. Symmetry Collapse Index (SCI) is S4 $0.889 \pm 0.002$, S2 $0.892 \pm 0.000$, S1 $0.984 \pm 0.002$; S4 versus S1 gives $p = 0.036$ and $r = -1.0$. Linear decoding error worsens: S4 $0.505 \pm 0.014$, S2 $0.448 \pm 0.014$, S1 $0.442 \pm 0.012$ ($p = 0.036$, $r = -1.0$). S4 compresses symmetry-related pairs by 11\% relative to random pairs. The spatial distribution of these errors is not uniform: decoding errors in S4 concentrate near the symmetry axes of the arena (Figure~\ref{fig:decode}), exactly where positions related by 90$^\circ$ rotation are hardest to distinguish. The dissociation between RA and SCI is notable: S4 and S2 show substantially different unit-level aliasing yet indistinguishable population compression ($p = 0.786$).

\begin{figure}[tbp]
\centering
\includegraphics[width=0.65\textwidth]{figures/fig_decoding_confusion.pdf}
\caption{Decoding confusion under symmetry. Errors concentrate near the symmetry axes --- exactly where 90$^\circ$ rotation makes positions hardest to distinguish.}
\label{fig:decode}
\end{figure}

\section*{Distance geometry shifts toward path-dependence under symmetry}

Symmetry transforms distance geometry. Distance-Topology Gap (DTG) is S4 $-0.052 \pm 0.003$, S2 $-0.011 \pm 0.003$, S1 $-0.020 \pm 0.004$ ($p = 0.036$, $r = 1.0$). All gaps are negative, indicating city-block dominance, but the gap widens in S4. Under $C_4$ symmetry, inter-quadrant Euclidean structure collapses because identical observations push those positions toward similar hidden states; path length within quadrants dominates the learned distance geometry.

\section*{C2 symmetry is encoded as group structure in the population code}

The population code encodes group structure, not all symmetries uniformly. A negative C2 Contrast (mean$_{\text{C2}}$ distance $<$ mean$_{\text{C4}}$ distance) indicates that positions related by 180\degree\ are more similar, consistent with C2 subgroup encoding. In S4, all quadrants are equally confusable, so inter-quadrant pairs compress identically and subgroup structure vanishes.

\begin{figure}[tbp]
\centering
\includegraphics[width=0.65\textwidth]{figures/fig_quadrant_distances.pdf}
\caption{Inter-quadrant neural distances. C2-related quadrants are more similar in S2; all pairs compress uniformly in S4.}
\label{fig:quadrant}
\end{figure}

\section*{Manifold geometry stays high-dimensional and folds along symmetry axes}

The network maintains high dimensionality despite symmetry. TwoNN intrinsic dimensionality is S4 $5.95 \pm 0.14$, S2 $5.78 \pm 0.23$, S1 $5.66 \pm 0.26$.

\begin{figure}[tbp]
\centering
\includegraphics[width=0.45\textwidth]{figures/fig_intrinsic_dim_training.pdf}
\hfill
\includegraphics[width=0.52\textwidth]{figures/fig_population_geometry.pdf}
\caption{Manifold geometry. Left: intrinsic dimensionality over training --- increases with symmetry order. Right: t-SNE visualizations showing local fragmentation in S4 and S2 while global structure is preserved.}
\label{fig:manifold}
\end{figure}

\begin{figure}[tbp]
\centering
\includegraphics[width=0.65\textwidth]{figures/fig_spectral_gap.pdf}
\caption{Spectral gap of the recurrent weight matrix $\mathbf{W}_h$. The dominant eigenvalue grows relative to the rest of the spectrum as symmetry increases --- but this predicts worse spatial maps, not better.}
\label{fig:spectral}
\end{figure}

\noindent The manifold fragments locally while preserving global structure---partial folding. The network cannot fold symmetric positions onto a single point without destroying within-quadrant structure for multi-step prediction, so it maintains high dimensionality while tolerating inter-quadrant ambiguity. The spectral gap of $\mathbf{W}_h$ increases with symmetry order (Figure~\ref{fig:spectral}), but higher spectral gap goes with worse spatial maps --- the opposite of what a stronger attractor would predict. Across all results: global degeneracy is ruled out, local aliasing scales with symmetry, population compression stays graded, distance geometry shifts toward path-dependence, the population code encodes group structure, and the manifold stays high-dimensional but fragments. Symmetry costs are real but bounded and compensable.

\section*{Heading Direction Input is Necessary for Global Map Stability}
\label{sec:hdablation}

The PAA analysis in the symmetry sweep suggested that HD input plays a
central role in preventing global orientational degeneracy. To test this
directly, we trained three variants of the pRNN in the S4 arena (the
hardest condition): HD\_FULL (the standard model), HD\_ABLATED (heading
input zeroed: $\mathbf{d}_t[1:] = 0$), and HD\_DEGRADED (heading set to
uniform: $\mathbf{d}_t[1:] = 0.25$). Speed was preserved in all
conditions.

Removing heading input caused map formation to fail completely. Spatial
sRSA dropped from $\sim$0.615 in HD\_FULL to
$\sim$0.260 in HD\_ABLATED (Figure~\ref{fig:hdablation}),
well below the map formation gate of $0.40$. HD\_DEGRADED, which provides
heading information but destroys its directional specificity by making all
four heading channels equally active, produced the same failure
(sRSA\,$\approx$\,0.258). This shows that it is
specifically the directional content of the HD signal, not just its
presence, that matters for map formation.

The mechanism is visible in the training loss. The final loss for ablated
conditions ($\sim$0.0190) was about $8\times$ higher than
for HD\_FULL ($\sim$0.0025), indicating that the network
without heading direction simply cannot predict future observations
accurately even after 80k training steps. Without knowing where the agent
is headed, the optimal strategy becomes to predict the mean observation ---
which carries no spatial information. Mean field coherence confirmed this:
HD\_ABLATED networks had a coherence of $\sim$58.96
(highly disordered activity), compared to $\sim$13.64 for
HD\_FULL.

Intrinsic manifold dimensionality was also lower in ablated networks
(ID\,$\approx$\,5.34 vs 6.09 for
HD\_FULL), suggesting that the loss of spatial structure is accompanied by
a genuine reduction in the representational degrees of freedom, not just
a degradation in spatial tuning. Together, these results confirm that HD
input is not a passive contextual signal but a necessary input for the
predictive learning objective to produce spatially organized representations.
This is consistent with the known role of head direction signals in
anchoring hippocampal spatial representations to a global reference
frame~\citep{lever2002}.

\begin{figure}[tbp]
\centering
\includegraphics[width=0.65\textwidth]{figures/fig_hd_ablation.pdf}
\caption{HD ablation results. Removing heading input causes map collapse (sRSA drops below the 0.40 formation gate). Degrading directional specificity produces the same failure.}
\label{fig:hdablation}
\end{figure}



\section*{Predictions the model generates for hippocampal biology}

The model predicts that multi-field place cell fraction should scale with symmetry group order. This follows directly from the RA result: existing hippocampal recording datasets from symmetric arenas can be stratified by symmetry order and the fraction of multi-field units measured. The RA result predicts that multi-field fraction scales as $|C_4| > |C_2| > |C_1|$, matching the symmetry groups: RA = 0.217, 0.114, 0.097 respectively. This prediction follows from the computational finding that high-RA units respond identically at rotationally equivalent positions and cannot disambiguate arena quadrant---precisely the dual-frequency structure of multi-field place cells recorded in symmetric arenas~\cite{muller1987, knierim1995}.

The model generates a specific prediction about head-direction lesion effects. Lesioning anterior thalamic nuclei (which provide HD input to hippocampus) should raise PAA gain above 0.05 in S4 arenas while RA remains at 0.217 or higher. This dissociation is the key prediction: global orientation becomes ambiguous between seeds (PAA breaks) but local aliasing persists (RA unchanged), because HD input stabilises global phase through transition statistics without affecting local observation equivalence. This dissociation is unique to the partial folding account.

The DTG result generates a testable prediction about population vector geometry. In symmetric arenas, population vector distances should correlate more strongly with path-length than Euclidean distance, with the correlation bias largest in $C_4$ arenas and intermediate in $C_2$ arenas. The DTG result quantifies this: city-block bias is $-0.052 \pm 0.003$ in S4 versus $-0.020 \pm 0.004$ in S1. This prediction is directly testable in existing hippocampal recording data.



\section*{Limitations and replication requirements}

Sample sizes are at the statistical minimum. The sample sizes are S4 $n=7$, S2 $n=5$, S1 $n=5$, giving minimum uncorrected $p$-values of 0.036 and minimum Bonferroni-corrected $p$-values of 0.107 (correction factor = 3 comparisons). All main comparisons sit at theoretical minimum uncorrected $p$ and maximum effect size ($r = -1.0$), indicating the effects are as large as statistically observable at these sample sizes. Replication requires $n \geq 12$ seeds per condition to achieve corrected significance ($p < 0.05$ after Bonferroni correction for 3 metrics).

The SCI and RA dissociation is unexplained. S4 and S2 show substantially different unit-level aliasing (RA: $0.217 \pm 0.003$ versus $0.114 \pm 0.004$) yet statistically indistinguishable population compression (SCI: $p = 0.786$). This dissociation suggests the population code redistributes aliasing across units rather than concentrating it through compensatory coding strategies. The mechanistic cause remains unidentified and requires analysis of learned weight matrices and representational geometry.

Distance geometry results confound representation with environment. The $18 \times 18$ MiniGrid imposes a city-block navigation metric by construction---agents must move along grid lines. The DTG result showing larger negative gaps under symmetry may partially reflect this environmental structure. Replication in a continuous environment such as RatInABox is required before the representational geometry interpretation is secure.

Spatially tuned unit fraction is lower than typical recording values. The fraction of spatially tuned units at EVS $> 0.1$ is S4 1.8\%, S2 3.7\%, S1 4.7\%, substantially lower than values typically reported in hippocampal recordings at more lenient thresholds (EVS $> 0.5$). The more conservative threshold reflects position-averaged evaluation which removes within-position variability. The ordering $S4 < S2 < S1$ is robust to threshold choice and the main conclusions rest on population-level metrics (RA, SCI, DTG, C2 Contrast) unaffected by this discrepancy.

The HD ablation used $n=2$ seeds and should be treated as a targeted validation rather than a fully-powered statistical test.



\section{Discussion}
\label{sec:discussion}

\subsection{Partial Folding as the Generic Response to Environmental Symmetry}

The central finding of this work is that the pRNN framework does not fail
under landmark symmetry. A network trained to
predict future observations in a C4-symmetric environment does not produce
a globally degenerate or random representation. Instead, it forms a map
--- but one that is geometrically compressed along the symmetry axes of
the environment. We call this partial folding: the network learns the true
spatial structure of the arena, but units develop activity patterns that
are correlated at rotationally equivalent positions, reducing the
precision of the spatial code without destroying it.

This is the right outcome from a computational standpoint. The task does
not require the network to distinguish between positions that are
genuinely indistinguishable from the agent's sensory perspective. Under
C4 landmark symmetry, positions separated by 90\degree~do produce the
same observations on average. The network's representational geometry
honestly reflects this ambiguity. The SR framework predicts exactly this
outcome: if two positions have similar expected future state distributions
(as they must when the arena looks the same from both), their SR entries
will be similar, and a network that learns the SR will represent them as
nearby in neural space~\citep{stachenfeld2017}. What we observe as partial
folding is the network faithfully encoding the transition structure of a
symmetric environment.

The spectral gap of $\mathbf{W}_h$ increases with symmetry order
(Figure~\ref{fig:spectral}), but negatively predicts spatial map quality.
This looks like a signature of a stronger, more stable attractor ---
but it predicts worse spatial representations across seeds. The attractor
has organized around the environmental symmetry group rather than the
arena's spatial layout. Attractor strength and attractor accuracy are not
the same thing, and symmetry can drive them apart.

\subsection{The Role of Heading Direction}

The HD ablation experiment clarifies a distinction that is easy to
conflate: the difference between global orientational degeneracy (the
network cannot tell which way it is facing) and local representational
aliasing (the network represents multiple positions similarly). HD input
prevents the first but not the second. Without heading direction, the
predictive objective becomes unachievable, and the network abandons the
attempt to form a spatial representation entirely. With heading direction,
a stable global orientation is maintained (PAA gain $< 0.05$), but the
local folding induced by the landmark structure persists.

This dissociation maps naturally onto what is known about the biological
system. Head direction cells in the anterior thalamus and entorhinal cortex
provide a consistent allocentric orientation signal that is thought to be
necessary for anchoring hippocampal place fields to a global reference
frame~\citep{lever2002}. Our results suggest that this anchoring function
is precisely what heading input provides to the predictive RNN: it resolves
the ambiguity about global orientation, but cannot resolve the local
ambiguity created by observation aliasing. The remaining
aliasing reflects environmental structure rather than a failure of the
network.

\subsection{Representational Quality and the Spectral Gap}

The spectral analysis reveals a negative correlation between representational quality (measured by sRSA) and the spectral gap of the recurrent weight matrix. The spectral gap --- the distance between the largest and second-largest eigenvalues --- reflects the stability and "mixing time" of the recurrent dynamics. In symmetric environments, we find that the spectral gap narrows as sRSA increases. This suggests that high-quality spatial maps in symmetric arenas do not emerge from "stronger" or more stable recurrent attractors, but rather from dynamics that are poised near a symmetry-breaking transition. This poise allows the network to maintain high-dimensional within-quadrant structure while tolerating the global ambiguity induced by the landmarks. The spectral gap thus serves as a dynamical signature of the "partial folding" we observe at the representational level.

\subsection{Limitations and Future Directions}

Several limitations of this work should be noted. The HD ablation
experiment used $n=2$ seeds and should be treated as a targeted
validation rather than a fully-powered statistical test.

More broadly, this work characterizes the representational geometry that
emerges from predictive learning but does not address whether this
geometry supports downstream computations. Place cells in the hippocampus
are read out by downstream regions for navigation and memory. Whether the
partial folding we observe degrades these downstream computations in
proportion to the RA increase, or whether readout mechanisms compensate
for the geometric distortion, is an important open question. Future work
extending this analysis to multi-environment generalization --- where
the network must maintain distinct maps for environments with different
symmetry structures --- will be needed to address this.






\bibliographystyle{plainnat}
\bibliography{references}

\appendix
\section{Implementation Details}
\label{app:implementation}

\subsection*{Architecture}
We use the pRNN\textsubscript{th} (rollout) architecture of
\citet{levenstein2024} with $N=500$ hidden units and NormReLU nonlinearity:
\begin{equation}
    \mathbf{h}_t = \text{NormReLU}\!\left(\mathbf{W}_h\mathbf{h}_{t-1} +
    \mathbf{W}_o\mathbf{o}_t + \mathbf{W}_a\mathbf{a}_t +
    \mathbf{W}_d\mathbf{d}_t\right)
\end{equation}
where $\text{NormReLU}(\mathbf{x}) = \text{ReLU}(\mathbf{x}) /
\|\text{ReLU}(\mathbf{x})\|_2$. The input $\mathbf{d}_t$ is a 5-dimensional
SpeedHD vector: $d_0 =$ locomotion speed, $d_{1:4} =$ 4-way one-hot
heading direction. The rollout depth is $k=5$ steps and the sequence
length is $T=200$. All architectural choices follow Levenstein et al.
except where noted.

\subsection*{Optimizer}
We use RMSProp with $\alpha=0.95$ and $\varepsilon=10^{-9}$. This follows
the optimizer choice of the original Levenstein et al. implementation.
Adam was not used because the per-parameter learning-rate scaling in this
RNN already performs much of Adam's adaptive normalization; using Adam in
addition risks over-adapting the gradients in a temporally correlated
sequence setting.

\subsection*{Training}
Each run was trained for 80,000 steps. Gradient accumulation over 8 steps
(ACCUM\_STEPS=8) was used in place of true batch size 8, which would have
exceeded available GPU memory. Gradient accumulation is mathematically
equivalent for this optimizer and architecture.

\subsection*{Evaluation}
Spatial sRSA was computed using Spearman correlation between the upper
triangles of the neural pairwise distance matrix (cosine) and the spatial
pairwise distance matrix (Euclidean), over a fixed evaluation trajectory
of 324 positions. The EVS threshold for classifying a unit as spatially
tuned was set to $0.1$ (confirmed from the \texttt{evs\_threshold} key
in evaluation outputs). All reported statistics use the values from the
final evaluation checkpoint (\texttt{eval.pkl}), not from intermediate
checkpoints, to ensure consistent trajectory length across seeds.
\end{document}